\documentclass[dvipdfmx, a4paper, 11pt, titlepage]{jsarticle}
\usepackage[dvipdfmx]{graphicx}
\usepackage{listings}
\usepackage{amsmath}
\usepackage{url}
\usepackage{booktabs}

\title{知能情報実験III（データマイニング班）\\競馬複勝予測}
\author{グループの学籍番号列挙 245725D, 245304F, 245731J, 235751J}
\date{提出日：2026年6月9日}
\begin{document}
\maketitle
\tableofcontents
\clearpage


\section{本実験向け補足}

\subsection{競馬複勝予測とは}
本グループでは、競馬における馬の複勝率予測を対象問題として設定した。複勝とは、出走した馬が3着以内に入るかどうかを予測する馬券の種類であり、単勝や馬連と比較して的中しやすい特徴を持つ。本テーマでは、過去の競馬データを用いて、各馬が複勝するか否かを機械学習によって分類することを目的とする。
競馬では、馬の年齢や体重、レース距離、天候、馬場状態など多くの要因がレース結果に影響すると考えられている。しかし、人間がこれらの膨大な情報を総合的に判断して予測を行うことは容易ではない。そこで、機械学習を活用することで、多数の説明変数から複勝に関係する特徴を学習し、客観的な予測を行うことが可能となる。

本テーマに取り組む意義は、実際の競馬予想への応用だけでなく、多数の特徴量を持つ実データに対して分類モデルを適用し、その性能や特徴量の重要度を分析できる点にある。これにより、機械学習を用いた予測手法の有効性を検証するとともに、データ分析の実践的な知識を深めることができる。

\subsection{実験計画（中間報告用）}
本グループでは4名で役割分担を行い、2名がデータの加工を担当し、残りの2名が機械学習モデルの実装および評価を担当する。

まず、1週目はデータ担当者が競馬の過去レースデータを収集し、欠損値や異常値の確認、不要な列の削除などの前処理を行う。また、複勝した馬を1、複勝しなかった馬を0とする教師データを作成する。一方、機械学習担当者は使用するライブラリやモデルの調査を行い、学習環境の構築およびベースラインモデルの作成を進める。

2週目は、データ担当者がカテゴリ変数のエンコーディングや特徴量の整理を行い、学習に利用できるデータセットを完成させる。機械学習担当者は完成したデータセットを用いて学習を実施し、正解率や適合率、再現率などの評価指標を算出する。また、特徴量重要度の分析を行い、予測に大きく影響する要素を調査する。

その後は、得られた結果を基にPDCAサイクルを回す。まず学習結果を分析し、予測精度が低い場合には原因を検討する。例えば、特徴量不足やデータの偏り、モデル選択の問題などを考察する。その分析結果を踏まえて、新たな特徴量の追加や不要特徴量の削除、ハイパーパラメータの調整、別モデルの適用などの改善策を実施する。改善後のモデルを再度評価し、性能向上の有無を確認する。

最終的には、複数のモデルの性能比較を行い、どのような特徴量が複勝予測に有効であったかを考察する。また、実験結果をまとめ、機械学習による競馬予測の有効性と課題について報告する予定である。



\subsection{実験目的}
本実験の目的は、競馬の過去データを用いて機械学習モデルを構築し、出走馬が複勝（3着以内に入着）するかどうかを予測できるかを検証することである。競馬のレース結果には、馬齢、馬体重、距離、天候、馬場状態など多くの要因が関係していると考えられるため、これらの特徴量から複勝の可否を予測する分類モデルを作成する。
また、複数の特徴量を用いた機械学習によってどの程度の予測性能が得られるのかを評価するとともに、どの特徴量が複勝予測に大きく影響しているのかを分析することも目的とする。さらに、異なる機械学習手法や前処理方法を比較し、それらが予測精度に与える影響を検証する。
これらの実験を通して、競馬における複勝予測の実現可能性を確認するとともに、実データを対象とした機械学習モデル構築の有効性について考察する。

\subsection{実験手順（予定）}

本実験では、Kaggleで公開されているJRA競馬データセットを用いて、出走馬が複勝（3着以内入着）するかどうかを予測する二値分類問題に取り組む予定である。

まず、データセットから競走馬の情報を抽出し、目的変数を作成する。目的変数は着順を基に作成し、3着以内を1、4着以下を0とする予定である。また、馬齢、馬体重、馬体重増減、距離、天候、馬場状態などの特徴量を説明変数として利用する。欠損値の処理やカテゴリ変数のエンコーディングなどの前処理を行い、機械学習で利用可能な形式へ変換する。

次に、データを訓練データとテストデータに分割し、訓練データを用いて機械学習モデルの学習を行う予定である。モデルとしてはロジスティック回帰やXGBoostなどを候補とし、性能を比較する。

評価には正解率（Accuracy）、F1スコアを用いる予定である。さらに、特徴量やハイパーパラメータを変更しながら性能を比較し、複勝予測に適したモデルを検討する。

最終的には、各モデルの評価結果を比較し、競馬の複勝予測に対する機械学習の有効性について考察する予定である。


\subsection{データセット}
\url{https://www.kaggle.com/datasets/takamotoki/jra-horse-racing-dataset?select=19860105-20210731_race_result.csv}

このデータセットは、JRA（日本中央競馬会）の競馬データを収録したものであり、1986年から2021年までのレース結果、オッズ、ラップタイム、コーナー通過順位などの情報を含む。データは競馬情報サイト「netkeiba.com」から収集・整形されたものである。
本研究では、その中の race\_result.csv を主に利用する。このファイルには各レースに出走した馬ごとの情報が記録されており、馬体重、馬体重増減、馬齢、競馬場、距離、天候、馬場状態、騎手など、多数の説明変数が含まれている。

\subsection{モデル（予定）}
\begin{itemize}
 \item ロジスティック回帰
 \item XGBoost
 \item RandomForest
 \end{itemize}


\subsection{パラメータ調整}
特になし

\section{実験}
\subsection{使用した特徴量}
\begin{table}[htbp]
  \centering
  \caption{基本特徴量}
  \label{tab:basic_features}
  \begin{tabular}{ll}
    \hline
    特徴量名 & 説明 \\
    \hline
    馬齢 & 競走馬の年齢 \\
    馬体重 & レース当日の馬体重 \\
    馬体重増減 & 前走からの馬体重の増減 \\
    斤量 & 負担重量 \\
    性別 & 牡・牝・セン \\
    距離 & レース距離(m) \\
    芝・ダート区分 & 芝コースかダートコースか \\
    枠番 & 出走時の枠番号 \\
    \hline
  \end{tabular}
\end{table}

\begin{table}[htbp]
  \centering
  \caption{レース条件に関する特徴量}
  \label{tab:race_condition_features}
  \begin{tabular}{ll}
    \hline
    特徴量名 & 説明 \\
    \hline
    競馬場名 & 開催競馬場 \\
    馬場状態 & 良・稍重・重・不良 \\
    天候 & レース当日の天候 \\
    右左回り・直線区分 & コースの回り方向 \\
    距離帯 & 距離を5区分にカテゴリ化(短距離・マイル・中距離・中長距離・長距離) \\
    レースクラス & 新馬・未勝利・1〜3勝クラス・オープンの6区分 \\
    レースグレード & G1・G2・G3・L・なしの5区分 \\
    休養日数 & 前走からの日数 \\
    \hline
  \end{tabular}
\end{table}

\begin{table}[htbp]
  \centering
  \caption{過去成績に基づく特徴量}
  \label{tab:past_performance_features}
  \begin{tabular}{ll}
    \hline
    特徴量名 & 説明 \\
    \hline
    過去3走・5走平均着順 & 直近の平均着順 \\
    過去3走・5走複勝率 & 直近の複勝(3着以内)率 \\
    過去3走・5走平均上り & 直近のラスト3F平均タイム \\
    騎手複勝率 & 騎手の通算複勝率 \\
    調教師複勝率 & 調教師の通算複勝率 \\
    芝・ダート別複勝率 & 芝/ダート別の複勝率 \\
    距離帯別複勝率 & 距離帯別の複勝率 \\
    競馬場別複勝率 & 競馬場別の複勝率 \\
    \hline
  \end{tabular}
\end{table}

すべての過去成績に基づく特徴量は、対象レースより過去のデータのみを用いて算出しており、
未来の情報が混入しないよう\texttt{shift}処理を施している。

\begin{table}[htbp]
  \centering
  \caption{同一レース内での相対化特徴量}
  \label{tab:relative_features}
  \begin{tabular}{ll}
    \hline
    特徴量名 & 説明 \\
    \hline
    順位率 & 同一レース内でのパーセンタイル順位(0に近いほど上位) \\
    標準化 & 同一レース内での標準化得点(Zスコア) \\
    \hline
  \end{tabular}
\end{table}
上記の過去成績特徴量および馬体重・斤量・馬齢・休養日数について、
同一レースに出走する馬同士で比較した相対値を追加で作成した。
\subsection{実験結果}

\subsubsection{実験1:ベースライン}

使用した特徴量:馬齢, 馬体重, 距離(m), 芝・ダート区分

\begin{table}[htbp]
  \centering
  \caption{実験1の結果}
  \label{tab:exp1}
  \begin{tabular}{lccccc}
    \toprule
    モデル & Accuracy & F1 & ROC-AUC & Precision & Recall \\
    \midrule
    XGBoost           & 0.4794 & 0.3544 & 0.5595 & 0.2427 & 0.6565 \\
    Random Forest     & 0.5012 & 0.3363 & 0.5408 & 0.2367 & 0.5804 \\
    ロジスティック回帰 & 0.4637 & 0.3544 & 0.5569 & 0.2401 & 0.6762 \\
    \bottomrule
  \end{tabular}
\end{table}

\subsubsection{実験2:市場予想(単勝オッズ・人気)の追加}

実験1の特徴量に加え、単勝オッズ・人気を追加した。

\begin{table}[htbp]
  \centering
  \caption{実験2の結果}
  \label{tab:exp2}
  \begin{tabular}{lccccc}
    \toprule
    モデル & Accuracy & F1 & ROC-AUC & Precision & Recall \\
    \midrule
    XGBoost           & 0.7390 & 0.5456 & 0.8100 & 0.4393 & 0.7196 \\
    Random Forest     & 0.7551 & 0.4730 & 0.7611 & 0.4449 & 0.5048 \\
    ロジスティック回帰 & 0.7101 & 0.5368 & 0.8068 & 0.4115 & 0.7719 \\
    \bottomrule
  \end{tabular}
\end{table}

\subsubsection{実験3:市場予想を除外した特徴量}

単勝オッズ・人気を除外し、代わりに斤量, 馬体重増減, 枠番, 休養日数, 性別,
競馬場名, 馬場状態, 天候, 右左回り・直線区分, 距離帯, レースクラス, レースグレード
を追加した。

\begin{table}[htbp]
  \centering
  \caption{実験3の結果}
  \label{tab:exp3}
  \begin{tabular}{lccccc}
    \toprule
    モデル & Accuracy & F1 & ROC-AUC & Precision & Recall \\
    \midrule
    XGBoost           & 0.5448 & 0.3724 & 0.6009 & 0.2660 & 0.6203 \\
    Random Forest     & 0.7145 & 0.2233 & 0.5655 & 0.2737 & 0.1885 \\
    ロジスティック回帰 & 0.4681 & 0.3660 & 0.5770 & 0.2471 & 0.7053 \\
    \bottomrule
  \end{tabular}
\end{table}

\subsubsection{実験4:過去成績特徴量の追加}

実験3の特徴量に加え、過去平均着順, 過去複勝率, 調教師複勝率, 騎手複勝率,
距離帯別複勝率, 過去平均上り, 芝・ダート別複勝率, 競馬場別複勝率を追加した。

\begin{table}[htbp]
  \centering
  \caption{実験4の結果}
  \label{tab:exp4}
  \begin{tabular}{lccccc}
    \toprule
    モデル & Accuracy & F1 & ROC-AUC & Precision & Recall \\
    \midrule
    XGBoost           & 0.7069 & 0.4806 & 0.7449 & 0.3913 & 0.6230 \\
    Random Forest     & 0.7712 & 0.4106 & 0.7356 & 0.4674 & 0.3662 \\
    ロジスティック回帰 & 0.7025 & 0.4667 & 0.7295 & 0.3827 & 0.5981 \\
    \bottomrule
  \end{tabular}
\end{table}

\subsubsection{実験5:同一レース内での相対化}

過去成績や馬体重などの特徴量について、同一レース内での順位率・標準化を追加した。

\begin{table}[htbp]
  \centering
  \caption{実験5の結果}
  \label{tab:exp5}
  \begin{tabular}{lccccc}
    \toprule
    モデル & Accuracy & F1 & ROC-AUC & Precision & Recall \\
    \midrule
    XGBoost           & 0.7110 & 0.4957 & 0.7612 & 0.3997 & 0.6524 \\
    Random Forest     & 0.7761 & 0.4358 & 0.7525 & 0.4825 & 0.3973 \\
    ロジスティック回帰 & 0.7028 & 0.4864 & 0.7492 & 0.3898 & 0.6466 \\
    \bottomrule
  \end{tabular}
\end{table}

\subsubsection{実験6:回収率による評価}

実験5のモデルを用いて、閾値方式による複勝回収率を評価した。
実験5のモデルを用いて、モデルの予測確率が一定の閾値を超えた場合にのみ
馬券を購入したと仮定し、複勝回収率を評価した。
購入方法として、以下の2種類を比較した。

\begin{itemize}
  \item \textbf{全頭方式}:閾値を超えた馬を、1レースにつき複数頭でも
        すべて購入する方式(該当馬が0頭のレースは購入しない)。
  \item \textbf{1位方式}:各レースで予測確率が最も高い馬のみを対象とし、
        その馬の予測確率が閾値を超えている場合にのみ購入する方式
        (1レースにつき購入は0頭または1頭のみ)。
\end{itemize}

いずれの方式も、1頭(1点)あたり100円を購入したものと仮定し、
以下の式で回収率を算出した。

\[
  \text{回収率} = \frac{\text{総払戻金}}{\text{総購入金額}} \times 100
\]

なお、閾値を高く設定するほど購入対象となる頭数(レース数)は減少するため、
回収率とあわせて実際に賭けた頭数・購入レース数も併記する。


\subsubsection*{XGBoost}

表~\ref{tab:exp6_xgb_all}に全頭方式、表~\ref{tab:exp6_xgb_top1}に1位方式の結果を示す。

\begin{table}[htbp]
  \centering
  \caption{実験6:XGBoost 回収率(閾値\&全頭方式)}
  \label{tab:exp6_xgb_all}
  \begin{tabular}{lccccc}
    \toprule
    閾値 & 0.5 & 0.6 & 0.7 & 0.8 & 0.9 \\
    \midrule
    回収率(\%)   & 80.05 & 80.43 & 81.17 & 82.84 & 85.65 \\
    賭けた頭数   & 475{,}808 & 317{,}233 & 177{,}845 & 68{,}606 & 8{,}332 \\
    \bottomrule
  \end{tabular}
\end{table}

\begin{table}[htbp]
  \centering
  \caption{実験6:XGBoost 回収率(閾値\&1位方式)}
  \label{tab:exp6_xgb_top1}
  \begin{tabular}{lccccc}
    \toprule
    閾値 & 0.5 & 0.6 & 0.7 & 0.8 & 0.9 \\
    \midrule
    回収率(\%)     & 81.16 & 81.24 & 81.47 & 82.70 & 85.61 \\
    購入レース数   & 97{,}535 & 96{,}248 & 86{,}245 & 51{,}705 & 8{,}081 \\
    \bottomrule
  \end{tabular}
\end{table}

\subsubsection*{Random Forest}

表~\ref{tab:exp6_rf_all}に全頭方式、表~\ref{tab:exp6_rf_top1}に1位方式の結果を示す。

\begin{table}[htbp]
  \centering
  \caption{実験6:Random Forest 回収率(閾値\&全頭方式)}
  \label{tab:exp6_rf_all}
  \begin{tabular}{lccccc}
    \toprule
    閾値 & 0.5 & 0.6 & 0.7 & 0.8 & 0.9 \\
    \midrule
    回収率(\%)   & 80.09 & 80.91 & 82.66 & 83.83 & 89.44 \\
    賭けた頭数   & 240{,}000 & 109{,}079 & 37{,}422 & 7{,}935 & 557 \\
    \bottomrule
  \end{tabular}
\end{table}

\begin{table}[htbp]
  \centering
  \caption{実験6:Random Forest 回収率(閾値\&1位方式)}
  \label{tab:exp6_rf_top1}
  \begin{tabular}{lccccc}
    \toprule
    閾値 & 0.5 & 0.6 & 0.7 & 0.8 & 0.9 \\
    \midrule
    回収率(\%)     & 80.28 & 80.90 & 82.40 & 83.70 & 89.39 \\
    購入レース数   & 93{,}657 & 68{,}701 & 31{,}154 & 7{,}527 & 555 \\
    \bottomrule
  \end{tabular}
\end{table}

\subsubsection*{ロジスティック回帰}

表~\ref{tab:exp6_logistic_all}に全頭方式、表~\ref{tab:exp6_logistic_top1}に1位方式の結果を示す。


\begin{table}[htbp]
  \centering
  \caption{実験6:ロジスティック回帰 回収率(閾値\&全頭方式)}
  \label{tab:exp6_logistic_all}
  \begin{tabular}{lccccc}
    \toprule
    閾値 & 0.5 & 0.6 & 0.7 & 0.8 & 0.9 \\
    \midrule
    回収率(\%)   & 79.33 & 79.31 & 80.09 & 81.19 & 84.39 \\
    賭けた頭数   & 483{,}452 & 299{,}497 & 160{,}847 & 62{,}786 & 8{,}634 \\
    \bottomrule
  \end{tabular}
\end{table}

\begin{table}[htbp]
  \centering
  \caption{実験6:ロジスティック回帰 回収率(閾値\&1位方式)}
  \label{tab:exp6_logistic_top1}
  \begin{tabular}{lccccc}
    \toprule
    閾値 & 0.5 & 0.6 & 0.7 & 0.8 & 0.9 \\
    \midrule
    回収率(\%)     & 79.81 & 79.91 & 80.36 & 81.18 & 84.43 \\
    購入レース数   & 97{,}572 & 95{,}383 & 84{,}069 & 51{,}026 & 8{,}539 \\
    \bottomrule
  \end{tabular}
\end{table}

\subsubsection{参考:単勝オッズ・人気を特徴量に含めたモデル}

実験5の特徴量に加えて単勝オッズおよび人気を追加したモデルについても、
同様の方法で複勝回収率を算出した。
なお、これらのモデルは人気(市場予想)そのものを特徴量として利用しているため、
モデルの予測が人気の情報を強く反映しやすい点に留意されたい。

\paragraph{XGBoost}

表~\ref{tab:exp6_xgb_odds_all}に全頭方式、表~\ref{tab:exp6_xgb_odds_top1}に1位方式の結果を示す。

\begin{table}[htbp]
  \centering
  \caption{参考:XGBoost(オッズ・人気を含む)回収率(閾値\&全頭方式)}
  \label{tab:exp6_xgb_odds_all}
  \begin{tabular}{lccccc}
    \toprule
    閾値 & 0.5 & 0.6 & 0.7 & 0.8 & 0.9 \\
    \midrule
    回収率(\%)   & 82.07 & 82.58 & 83.70 & 85.64 & 88.61 \\
    賭けた頭数   & 473{,}021 & 343{,}980 & 219{,}883 & 106{,}985 & 25{,}104 \\
    \bottomrule
  \end{tabular}
\end{table}

\begin{table}[htbp]
  \centering
  \caption{参考:XGBoost(オッズ・人気を含む)回収率(閾値\&1位方式)}
  \label{tab:exp6_xgb_odds_top1}
  \begin{tabular}{lccccc}
    \toprule
    閾値 & 0.5 & 0.6 & 0.7 & 0.8 & 0.9 \\
    \midrule
    回収率(\%)     & 84.15 & 84.15 & 84.21 & 85.31 & 88.66 \\
    購入レース数   & 97{,}595 & 97{,}571 & 96{,}330 & 77{,}999 & 24{,}759 \\
    \bottomrule
  \end{tabular}
\end{table}

\paragraph{Random Forest}

表~\ref{tab:exp6_rf_odds_all}に全頭方式、表~\ref{tab:exp6_rf_odds_top1}に1位方式の結果を示す。

\begin{table}[htbp]
  \centering
  \caption{参考:Random Forest(オッズ・人気を含む)回収率(閾値\&全頭方式)}
  \label{tab:exp6_rf_odds_all}
  \begin{tabular}{lccccc}
    \toprule
    閾値 & 0.5 & 0.6 & 0.7 & 0.8 & 0.9 \\
    \midrule
    回収率(\%)   & 82.33 & 83.66 & 85.19 & 86.90 & 90.18 \\
    賭けた頭数   & 283{,}057 & 164{,}133 & 76{,}504 & 23{,}721 & 2{,}754 \\
    \bottomrule
  \end{tabular}
\end{table}

\begin{table}[htbp]
  \centering
  \caption{参考:Random Forest(オッズ・人気を含む)回収率(閾値\&1位方式)}
  \label{tab:exp6_rf_odds_top1}
  \begin{tabular}{lccccc}
    \toprule
    閾値 & 0.5 & 0.6 & 0.7 & 0.8 & 0.9 \\
    \midrule
    回収率(\%)     & 83.63 & 83.87 & 85.01 & 86.89 & 90.10 \\
    購入レース数   & 97{,}403 & 90{,}379 & 60{,}840 & 22{,}644 & 2{,}752 \\
    \bottomrule
  \end{tabular}
\end{table}

\paragraph{ロジスティック回帰}

表~\ref{tab:exp6_logistic_odds_all}に全頭方式、表~\ref{tab:exp6_logistic_odds_top1}に1位方式の結果を示す。

\begin{table}[htbp]
  \centering
  \caption{参考:ロジスティック回帰(オッズ・人気を含む)回収率(閾値\&全頭方式)}
  \label{tab:exp6_logistic_odds_all}
  \begin{tabular}{lccccc}
    \toprule
    閾値 & 0.5 & 0.6 & 0.7 & 0.8 & 0.9 \\
    \midrule
    回収率(\%)   & 80.63 & 81.17 & 82.16 & 84.06 & 91.43 \\
    賭けた頭数   & 528{,}115 & 393{,}795 & 244{,}212 & 74{,}233 & 91 \\
    \bottomrule
  \end{tabular}
\end{table}

\begin{table}[htbp]
  \centering
  \caption{参考:ロジスティック回帰(オッズ・人気を含む)回収率(閾値\&1位方式)}
  \label{tab:exp6_logistic_odds_top1}
  \begin{tabular}{lccccc}
    \toprule
    閾値 & 0.5 & 0.6 & 0.7 & 0.8 & 0.9 \\
    \midrule
    回収率(\%)     & 83.09 & 83.09 & 83.09 & 84.03 & 91.43 \\
    購入レース数   & 97{,}595 & 97{,}595 & 97{,}559 & 66{,}058 & 91 \\
    \bottomrule
  \end{tabular}
\end{table}


\end{document}